% Projeto de Conclusão de Curso
% Aluno: xpto
%----------------------------
\documentclass[
	12pt,	
    oneside,
    a4paper,
    brazil  
]{abntex2}

% ----------------------------
% Pacotes para acentuação
% ----------------------------
    \usepackage[T1]{fontenc}
    \usepackage{lmodern}
    \usepackage{sty/sbc-template}
    \usepackage[utf8]{inputenc}
    \usepackage{adjustbox}    % Escala/resize
    \usepackage{longtable}    % Multi-página
    \usepackage{rotating}     % Paisagem
    \usepackage{booktabs}

%--------------------------------------------
% Variáveis para preencher com os seus dados
%--------------------------------------------
    \title{Desafios da modernização de sistemas legados com agentes de IA no desenvolvimento de software}
    \subtitle{}
    \author{Anna Beatriz Resende Oki Corrêa}
    \cidade{Brasília}
    \ano{2026}
    \coordenacao{Coordenação do Curso de Sistemas de Informação}
    \titulacao{ Sistemas de Informação}
    \orientador{Prof. Dr. Xxxx}
%--------------------------------------------


\renewcommand{\listfigurename}{}
\addto\captionsbrazil{\renewcommand{\tablename}{Quadro}}
\begin{document} 
\pagestyle{empty}
\maketitle
\newpage


% -------------------------------
% Registro Bibliográfico
% Faça as modificações nos campos que estão indicados com ??
% -------------------------------
\section*{}
\vspace*{80pt}
\begin{center}
\begin{tcolorbox}[
    enhanced, % Permite a personalização avançada
    sharp corners, % Torna os cantos afiados
    width=1.0\textwidth, % Largura da caixa
    colback=white, % Cor de fundo
    colframe=black, % Cor da primeira borda
    boxrule=3pt, % Largura da primeira borda
    boxsep=20pt, % Espaçamento entre o texto e as bordas
    frame style={draw=black, double distance=4pt, double, ultra thick}, % Estilo da segunda borda
    ]
    ?Sobrenome?, ?Nome?. 
    \hspace{1cm} Título : \titulo / ?author?. -- Brasília, 2026. 
    \hspace{1cm} \pageref{LastPage} f. \\[2mm]
   	Orientador: ?orientador?. 
    \hspace{1cm}Trabalho de Conclusão de Curso (Graduação -- Sistemas de Informação) -- Centro Universitário do Distrito Federal -- UDF. Coordenação do Curso de Sistemas de Informação, Brasília, DF, 2026. \\[2mm]
   	1. ?Assunto?. 2. ?Assunto?. 3. ?Assunto?. I. ?Título?. \\[2mm]
   	\begin{flushright}
    CDU: ?consultar na biblioteca?
    \end{flushright}
\end{tcolorbox}
\end{center}


% -------------------------------
% Assinaturas
% Faça as modificações nos campos que estão indicados com ??
% -------------------------------
\newpage
\section*{}

\vspace*{10pt}
\begin{center}
?Titulo?
\vspace*{60pt}
\begin{adjustwidth}{6cm}{-1cm}
Trabalho de conclusão de curso apresentado à ?Coordenação?, do Centro Universitário do Distrito Federal -- UDF, como requisito parcial para obtenção do grau de bacharel / tecnólogo em ?Curso?. \\[2mm]

Orientador: ?orientador?. 
\end{adjustwidth}
\vspace*{40pt}
Brasília, \underline{\hspace{2cm}} de \underline{\hspace{4cm}} de 2026.

\vspace*{60pt}
\begin{tabular}{p{7cm}}
    \hline
    \centering NOME DO EXAMINADOR \\
    \centering Titulação \\
    \centering Instituição a qual é filiado
\end{tabular}

\vspace*{30pt}
\begin{tabular}{p{7cm}}
    \hline
    \centering NOME DO EXAMINADOR \\
    \centering Titulação \\
    \centering Instituição a qual é filiado
\end{tabular}

\vspace*{30pt}
\begin{tabular}{p{7cm}}
    \hline
    \centering NOME DO EXAMINADOR \\
    \centering Titulação \\
    \centering Instituição a qual é filiado
\end{tabular}
\end{center}

\vspace*{40pt}
Nota: \underline{\hspace{2cm}}


% -------------------------------
% Dedicatórias
% -------------------------------
\newpage
\begin{dedicatoria2}{}
 Dedico este trabalho à minha família, pelo apoio, incentivo e compreensão durante toda a minha trajetória acadêmica.
 \end{dedicatoria2}

% -------------------------------
% Agradecimentos
% -------------------------------
\newpage
\begin{agradecimento}{Agradecimentos}
Agradeço, primeiramente, a Deus, por mais uma conquista; ao meu orientador, pela dedicação e correções; e aos bibliotecários pelo suporte em todas as pesquisas.

\end{agradecimento}


% -------------------------------
% Frase motivacional
% -------------------------------
\newpage
\begin{adjustwidth}{7cm}{0cm}
\vspace*{\fill}
\begin{flushright}
"A educação é a arma mais poderosa que você pode usar para mudar o mundo." \\
\vspace*{10pt}
Nelson Mandela
\end{flushright}
\end{adjustwidth}


% -------------------------------
% Resumo
% -------------------------------
\newpage
\begin{resumo2}
Este trabalho trata da atualização de sistemas legados para novas tecnologias com o uso de agentes de IA para desenvolvimento. O objetivo foi identificar os desafios desse uso, do ponto de vista do desenvolvimento. Para isso, foi feita uma revisão de estudos recentes sobre o tema, e os desafios encontrados foram organizados em categorias. Depois, foi feito um estudo de caso com um sistema real, feito em Power Apps. Com a ajuda de um agente de IA, quatro regras desse sistema foram convertidas para Python. Os resultados do código novo foram comparados, por testes, com os de um app de teste do Power Apps, que rodou as fórmulas originais com dados inventados. Nos casos testados, os resultados foram iguais. Mesmo assim, a leitura das fórmulas mostrou comportamentos escondidos nas regras, como uma regra que não ignora os registros excluídos e uma mesma conta feita de dois jeitos. Neste estudo, o desafio mais observado foi o das regras escondidas. Como foram poucas regras e os dados foram inventados, esses resultados valem só para o caso estudado.

\vspace*{20pt}
\textbf{Palavras-chave:} Sistemas legados. Modernização de software. Agentes de IA. Desenvolvimento de software.
\end{resumo2}


% -------------------------------
% Abstract - Resumo traduzido para o Inglês
% -------------------------------
\newpage
\begin{abstract}
This work addresses the update of legacy systems to new technologies using AI development agents. The objective was to identify the challenges of this use from a development perspective. To do so, a review of recent studies on the subject was carried out, and the challenges found were organized into categories. Then, a case study was carried out with a real system built in Power Apps. With the help of an AI agent, four rules of this system were converted to Python. The results of the new code were compared, through tests, with those of a Power Apps test app that ran the original formulas with made-up data. In the tested cases, the results were the same. Even so, reading the formulas revealed hidden behaviors in the rules, such as a rule that does not ignore deleted records and the same calculation done in two different ways. In this study, the most observed challenge was hidden rules. Since only a few rules were converted and the data was made up, these results apply only to the case studied.

\vspace*{20pt}
\textbf{Keywords:} Legacy systems. Software modernization. AI agents. Software development.
\end{abstract}


% -------------------------------
% Lista de Figuras
% -------------------------------
\newpage
\listoffigures*

% -------------------------------
% Lista de Tabelas
% -------------------------------



% -------------------------------
% Sumario
% -------------------------------
\tableofcontents*
%\cleardoublepage

 
% -------------------------------
% INICIO DO TRABALHO - CAPÍTULOS
% -------------------------------
\textual 
\chapter{Introdução}
Sistemas legados são sistemas que continuam em uso nas organizações, mesmo quando suas tecnologias ou ferramentas já não acompanham as práticas atuais. Isso não depende só da idade do sistema. Depende também da importância dele para os processos da organização e da dificuldade de mantê-lo, evoluí-lo e integrá-lo a novas soluções.

Muitas vezes, esses sistemas guardam regras de negócio criadas e alteradas ao longo de anos. Parte desse conhecimento pode estar só no código, principalmente quando a documentação é incompleta ou não acompanha as mudanças. Atualizar um sistema legado para tecnologias mais novas é o que se chama de modernização. Ela não é só reproduzir as funcionalidades em uma tecnologia recente: também é preciso entender e preservar as regras importantes para a organização. O sistema novo precisa continuar fazendo o que o antigo fazia. Por isso, esse trabalho exige análise e validação cuidadosas.

Estudos recentes trazem exemplos desse cenário. \citeonline{alves2026legacy} trabalharam com um sistema de gestão de shopping centers escrito em Visual Basic 6, uma linguagem antiga. O sistema estava em produção havia mais de vinte anos. Ele era organizado em um núcleo e em módulos secundários, acrescentados conforme o pedido de cada cliente da empresa dona do sistema. Os autores escolheram dois módulos de naturezas diferentes. O de Medições é voltado a cálculos numéricos, como consumo e tarifas. O de Compras depende de transações e de fluxos de aprovação. Os dois têm cerca de 299 mil e 224 mil linhas de código e são usados continuamente pelos clientes da empresa. Eles foram migrados para C\# (.NET 10), mantendo o mesmo banco de dados. A migração foi feita por um agente de IA, em uma única geração e sem feedback humano entre as tentativas. Ela envolveu 12 funcionalidades, seis de cada módulo: 6 simples, 4 médias e 2 complexas. O objetivo era que o código novo se comportasse como o original.

O segundo exemplo vem da IBM. De acordo com \citeonline{chakaravarthy2026cobol}, sistemas empresariais modernos muitas vezes dependem de aplicações antigas em COBOL, uma linguagem criada há várias décadas. A estrutura rígida e monolítica e as variáveis globais dessas aplicações dificultam a manutenção e a adoção da nuvem. Para lidar com isso, os autores propõem traduzir o COBOL para Java combinando a análise do código com um modelo de linguagem. A análise infere as classes e as assinaturas dos métodos em Java a partir dos dados e do fluxo de controle do programa em COBOL. Essas informações guiam o modelo na tradução da lógica dos programas. A solução faz parte de um produto da IBM, o watsonx Code Assistant for Z. Os autores a avaliaram em 20 programas de um conjunto chamado GenAPP e em 10 aplicações de clientes. Segundo eles, ela já é usada em vários ambientes de clientes da IBM.

Esses estudos também relatam dificuldades para o desenvolvedor. No estudo de \citeonline{alves2026legacy}, as regras espalhadas em vários arquivos e as regras escondidas em parâmetros opcionais e em configurações globais foram as mais difíceis de migrar. Outro estudo avaliou o uso de modelos de linguagem para documentar código em linguagens antigas, como MUMPS e assembly \cite{diggs2024leveraging}, para ajudar a entender esse tipo de código.

Uma forma de apoiar esse trabalho é usar agentes de IA para desenvolvimento. Eles são ferramentas que leem o código de um projeto, editam arquivos e rodam comandos e testes. Estudos recentes avaliaram esses agentes em tarefas de modernização \cite{li2026repomod, liu2026migrationbench, may2025freshbrew}.

Os resultados desses estudos relatam limites. No estudo de \citeonline{li2026repomod}, os agentes passaram em 91,3\% dos testes em projetos pequenos, mas em apenas 15,3\% nos projetos grandes. No estudo de \citeonline{alves2026legacy}, o agente deixou de migrar parte das regras do sistema, e as funcionalidades mais complexas foram as mais afetadas. Por isso, o desenvolvedor precisa saber quais desafios podem aparecer quando um agente de IA atualiza um sistema legado e como lidar com eles.

Este tema é relevante porque sistemas antigos precisam ser atualizados, mas esse trabalho é lento e arriscado. Um agente de IA pode acelerá-lo, mas a literatura relata desafios, como regras que o agente não percebe, atalhos que ele toma e resultados difíceis de verificar. Nos estudos que lemos, a avaliação usa conjuntos de testes públicos ou sistemas de empresas, e encontramos pouco sobre o uso do agente no dia a dia do desenvolvedor, em um sistema real e menor. Este trabalho busca ajudar o desenvolvedor a saber o que esperar e como se proteger. Para isso, faz uma revisão da literatura e faz um estudo de caso em um sistema real, com avaliação por testes.

Este trabalho busca responder à seguinte pergunta:

\begin{citacao}
Quais são os desafios para converter e atualizar sistemas legados para novas tecnologias por meio de agentes de IA, do ponto de vista do desenvolvimento?
\end{citacao}

Este trabalho está organizado em seis capítulos, contando esta introdução. O Capítulo 2 apresenta os objetivos. O Capítulo 3 reúne o que a literatura relata sobre sistemas legados, agentes de IA, verificação do código e os desafios desse uso. Com base nisso, o Capítulo 4 descreve como o estudo de caso foi feito, e o Capítulo 5 apresenta os resultados e as considerações finais. O Capítulo 6 traz a conclusão.


\chapter{Objetivos}
\label{cap:objetivos}
A introdução mostrou que os agentes de IA já são usados para modernizar sistemas legados, mas que se sabe pouco sobre como é usá-los no dia a dia de quem desenvolve, em um sistema real e de tamanho pequeno. Este capítulo diz o que o trabalho quer alcançar para tratar dessa lacuna e mostra em qual capítulo cada parte é feita.

O objetivo geral é identificar os desafios para converter e atualizar sistemas legados para novas tecnologias por meio de agentes de IA, do ponto de vista do desenvolvimento. Aqui, o ponto de vista do desenvolvimento é o de quem escreve, mantém e confere o código, e não o de quem administra ou usa o sistema. Um desafio, neste trabalho, é uma dificuldade que aparece quando o agente faz esse serviço e que o desenvolvedor precisa perceber e tratar, como uma regra do sistema que o agente não percebeu ou um resultado difícil de conferir. Esse objetivo responde à pergunta apresentada na introdução.

Para alcançá-lo, o trabalho olha para o assunto de dois lados: o que a literatura já relata e o que acontece na prática. Por isso, o objetivo geral é dividido em quatro etapas, que são os objetivos específicos, e cada uma é feita em um capítulo.

A primeira etapa é reunir o que estudos recentes relatam e organizar os desafios encontrados. O Capítulo~3, de fundamentação teórica, explica como esses estudos foram buscados e escolhidos e depois apresenta o assunto em quatro partes: como a modernização aparece nos estudos, como os agentes de IA são usados nessa tarefa, como o resultado é verificado e quais desafios a literatura relata. Esses desafios são reunidos em um quadro, que serve de lista do que deve ser observado no estudo de caso e que também mostra o que ainda falta ser estudado.

A segunda etapa é definir como o estudo de caso é feito. O Capítulo~4 explica o tipo de pesquisa e apresenta o sistema estudado, feito no Microsoft Power Apps: para que ele serve, como é organizado e quais regras dele foram escolhidas, de complexidades diferentes. Depois mostra que os dados usados são inventados, como os resultados do sistema original foram obtidos em um app de teste, como o agente de IA foi usado e como o resultado do código novo é comparado com o esperado. Assim, o leitor entende o que foi feito e consegue repetir o método.

A terceira etapa é converter essas regras para Python, com a ajuda de um agente de IA, e conferir o resultado por testes automáticos, que são pequenos programas que comparam o que o código devolve com o que se espera. O Capítulo~5 mostra o que foi convertido e o resultado dos testes.

A quarta etapa é observar os problemas que apareceram nessa conversão, ligá-los, quando possível, aos desafios do quadro da literatura e sugerir cuidados ao desenvolvedor. Esta parte também está no Capítulo~5, junto com os limites do estudo. O Capítulo~6 fecha o trabalho: responde à pergunta da introdução, retoma os limites e aponta o que pode ser feito em trabalhos futuros.

Para que o trabalho possa ser acompanhado de perto, ele faz algumas escolhas de escopo. O foco é o que o desenvolvedor encontra ao pedir a um agente de IA que atualize um sistema, e questões de gestão, de uso pelos usuários e de custo ficam de fora. Em vez de converter um sistema inteiro, o estudo de caso trata de quatro regras de um sistema feito no Microsoft Power Apps, de complexidades diferentes, que foram convertidas para Python. Também foi usado um só agente de IA, o que permite acompanhar o processo de perto, mas não permite comparar agentes entre si. A conferência do resultado é feita por testes automáticos, e a qualidade do código além disso, como desempenho ou organização, não é avaliada. Os limites que essas escolhas trazem para os resultados são discutidos no Capítulo~5.

A contribuição do trabalho é de três tipos. Primeiro, ele organiza em um quadro os desafios que a literatura relata, o que ajuda a enxergar, de uma vez só, o que pode dar errado nesse tipo de trabalho. Segundo, mostra em um caso pequeno e concreto como alguns desses desafios aparecem na prática e quais não puderam ser observados. Terceiro, deixa um método simples, com dados inventados, um app de teste e testes automáticos, que outro desenvolvedor pode repetir em um sistema seu.


\chapter{Fundamentação Teórica}
\label{cap:fundamentacao}
\section{Sistemas legados e modernização}

Um sistema legado é um sistema antigo que ainda está em uso. Ele pode estar escrito em uma linguagem pouco usada hoje, ter uma estrutura difícil de mudar ou depender de bibliotecas sem suporte. Mesmo assim, ele costuma ser importante, porque guarda regras de negócio que a organização usa todos os dias.

Um exemplo é o sistema estudado por \citeonline{alves2026legacy}. Era um sistema de gestão em Visual Basic 6, em produção havia mais de vinte anos. Outro exemplo são os programas em COBOL, que a IBM migra para Java \cite{chakaravarthy2026cobol}.

A modernização é o trabalho de levar esse sistema para uma tecnologia mais nova, sem perder o que ele já faz. Ela pode ser de dois tipos:

\begin{itemize}
    \item \textbf{conversão de linguagem}, por exemplo de Visual Basic 6 para C\#;
    \item \textbf{atualização de versão}, por exemplo de Java 8 para Java 17 \cite{liu2026migrationbench, may2025freshbrew}.
\end{itemize}

Uma dificuldade comum é a falta de documentação. Muitas vezes, o código é a única fonte do que o sistema faz. \citeonline{diggs2024leveraging} testaram modelos de linguagem para documentar código legado em MUMPS e em assembly. Entender o código antigo é o primeiro passo da modernização.

\section{Agentes de IA no desenvolvimento}

A seção anterior mostrou que modernizar um sistema exige entender o que ele faz e manter suas regras. Um agente de IA para desenvolvimento, como definido na introdução, lê os arquivos do projeto, altera o código, executa comandos e roda testes. Esta seção mostra como isso funciona em uma modernização e como os estudos avaliaram esses agentes.

O estudo de \citeonline{alves2026legacy} mostra bem como o agente é preparado para a tarefa. Os autores usaram o Claude Code, versão 2.1.96, com o modelo Opus 4.6. Para cada funcionalidade, criaram uma pasta de trabalho separada, com três arquivos de instruções, que são textos onde se escrevem o objetivo, as regras e as convenções que o agente deve seguir. O código antigo ficava em uma pasta somente de leitura, e o projeto novo, em C\#, partia de um modelo pronto.

O pedido era sempre o mesmo, padronizado: migrar a funcionalidade da pasta, lendo as instruções primeiro e usando o código antigo só como referência. O agente trabalhou em uma única passada, e o primeiro autor supervisionou as sessões e verificou o resultado.

Outros estudos avaliam os agentes com \textit{benchmarks}, que são conjuntos de projetos e testes preparados para comparar agentes nas mesmas condições. \citeonline{may2025freshbrew} criaram o FreshBrew, com 228 projetos Java públicos que compilam com o Java 8, falham no Java 17 e têm alta cobertura de testes, isto é, uma parte grande do código é conferida por testes. \citeonline{liu2026migrationbench} criaram o MigrationBench, com 5.102 repositórios e um subconjunto de 300, também para a migração do Java 8 para o 17. \citeonline{li2026repomod} criaram o RepoMod-Bench, com 21 repositórios em 8 linguagens, cerca de 1,6 milhão de linhas de código e 11.616 testes, para conversões entre pares de linguagens, como de Python para Java.

Os \textit{benchmarks} permitem comparar agentes de forma justa, mas usam projetos públicos e preparados para teste. Isso mostra bem o desempenho do agente, mas não mostra como é usá-lo em um sistema real do dia a dia, ponto que volta na Seção~\ref{sec:desafios-literatura}. O processo da IBM, visto na seção anterior, completa o quadro: em vez de um agente que conduz tudo, a análise do código guia o modelo de linguagem. Nos dois casos, fica uma pergunta: quem decide se o código novo está certo? É o tema da próxima seção.

\section{Conversão de código e verificação}

Converter código não é só trocar a sintaxe. O código novo precisa fazer a mesma coisa que o antigo. Isso é chamado de equivalência funcional.

Por isso, \citeonline{pan2024lost} consideram uma tradução correta só quando o código compila, roda e passa nos testes. Compilar não basta, porque um código pode compilar e mesmo assim se comportar de forma diferente.

Os testes também têm limites. Se o sistema tem poucos testes, erros passam sem ser vistos. Por isso, o \citeonline{may2025freshbrew} escolheu projetos com alta cobertura de testes e exige que a cobertura não caia mais de 5 pontos percentuais depois da migração. Já o \citeonline{li2026repomod} esconde os testes do agente, para que ele não os use como atalho.

Outra forma de avaliar é pedir para outro modelo de linguagem julgar o resultado. A IBM usou esse método para medir a qualidade do código Java gerado \cite{chakaravarthy2026cobol}. Mas, no estudo de \citeonline{diggs2024leveraging}, as métricas automáticas tiveram pouca relação com a nota dada por pessoas, e por isso os autores indicam a avaliação humana. Neste trabalho, a verificação principal será feita com testes.

\section{Desafios encontrados na literatura}
\label{sec:desafios-literatura}

A verificação mostra se o código novo está certo, mas os estudos relatam outros desafios além dela. Eles estão resumidos no Quadro~\ref{tab:desafios} e são explicados em seguida.

\begin{table}[htb]
\centering
\caption{Desafios da modernização com agentes de IA}
\label{tab:desafios}
\small
\begin{tabular}{p{3.2cm}p{6.6cm}p{3cm}}
\toprule
\textbf{Desafio} & \textbf{O que os estudos relatam} & \textbf{Fonte} \\
\midrule
Tamanho do sistema & Passou em 91,3\% dos testes com menos de 10 mil linhas e em 15,3\% com mais de 50 mil & \citeonline{li2026repomod} \\
Complexidade das funcionalidades & Equivalência de 92\% (baixa), 81\% (média) e 47\% (alta) & \citeonline{alves2026legacy} \\
Regras escondidas & O agente perdeu regras implícitas e regras espalhadas em vários arquivos & \citeonline{alves2026legacy} \\
Atalhos do agente & Excluiu testes que falhavam e, em 5 de 10 projetos, executou o programa original em vez de traduzi-lo & \citeonline{may2025freshbrew}; \citeonline{li2026repomod} \\
Linguagens antigas & Bons resultados em MUMPS, notas ruins em assembly, e alucinações em ambos & \citeonline{diggs2024leveraging} \\
Custo & Em média, 3,81 milhões de tokens e US\$ 4,19 por migração & \citeonline{alves2026legacy} \\
\bottomrule
\end{tabular}
\\[2pt]
\footnotesize Fonte: elaborado pela autora, com base nos trabalhos citados.
\end{table}

\subsection*{Tamanho e complexidade}

Os resultados costumam piorar quando o sistema é maior ou mais complexo. Quatro estudos mostram isso:

\begin{itemize}
    \item \citeonline{li2026repomod}: os agentes passaram em 91,3\% dos testes em projetos com menos de 10 mil linhas de código e em 15,3\% nos projetos com mais de 50 mil.
    \item \citeonline{alves2026legacy}: a equivalência com o sistema original foi de 92\% nas funcionalidades de complexidade baixa, 81\% nas médias e 47\% nas altas.
    \item \citeonline{liu2026migrationbench}: com o modelo Claude 4.5 Sonnet, houve 71,67\% de sucesso na migração mínima do Java 8 para o 17 e 53,33\% na máxima, em um subconjunto de 300 repositórios.
    \item \citeonline{may2025freshbrew}: o melhor modelo migrou 52,3\% dos projetos para o Java 17.
\end{itemize}

Os estudos usam medidas diferentes, por isso os números não devem ser comparados entre si. O que eles mostram em comum é que, nos casos relatados, o resultado cai em tarefas maiores ou mais difíceis e que nenhum agente acertou tudo.

\subsection*{Regras que o agente não vê}

Outro desafio é o agente não perceber regras do sistema antigo. \citeonline{alves2026legacy} mapearam 331 instruções nas 12 funcionalidades, sendo 160 regras de negócio e 171 operações de banco de dados. O agente deixou de migrar 100 delas. A maior causa, com 42 casos, foram módulos vizinhos que ele não encontrou. Depois vieram 22 operações de banco de dados e 20 regras que ficaram de fora.

Os autores também observaram que o agente tratou melhor as regras escritas de forma explícita do que as implícitas, como as que ficam escondidas em parâmetros opcionais ou em configurações dinâmicas. Isso sugere que o agente migra bem o que consegue ver, o que é um risco em sistemas com regras espalhadas.

\subsection*{Atalhos e falsa impressão de sucesso}

O agente pode parecer ter acertado sem ter acertado. \citeonline{may2025freshbrew} alertam que um agente poderia conseguir uma migração só na aparência, apagando os testes que falham no Java novo. Eles mostram casos em que um modelo tirou dois testes que falhavam do arquivo de construção do projeto e outro fez um teste ser pulado em Java moderno.

\citeonline{li2026repomod} viram algo parecido em um experimento inicial, em uma versão do teste sem regras contra atalhos. Um agente chegou a 99,5\% de testes aprovados, mas em 5 de 10 projetos usou códigos que apenas executavam o programa original. Depois de acrescentar essas regras, os autores relatam que esse comportamento deixou de aparecer. Por isso, o desenvolvedor deve proteger os testes e conferir o que o agente de fato fez.

\subsection*{Linguagens antigas e pouca documentação}

Em código legado, os modelos também podem errar de outro jeito. \citeonline{diggs2024leveraging} observaram que eles tendem a inventar informações, o que se chama de alucinação. Os comentários gerados para MUMPS foram avaliados por especialistas como bem melhores, em média, do que os de assembly, que tiveram qualidade baixa no geral. Os autores também lembram que outro estudo apontou a falta de documentação como obstáculo para saber se um código traduzido está correto.

\subsection*{Custo}

O custo depende do tamanho da tarefa. Em \citeonline{alves2026legacy}, cada migração usou, em média, 3,81 milhões de \textit{tokens}, que são as unidades em que os modelos contam o texto que leem e escrevem, e custou US\$ 4,19. As funcionalidades de alta complexidade usaram, em média, 9,09 milhões de \textit{tokens}. Quase todos os \textit{tokens}, 99,6\%, são do que o agente recebe, porque ele relê o código antigo a cada ciclo de trabalho. Por isso, os autores concluem que o custo de uma migração depende, em boa parte, do volume do que o agente recebe.

\subsection*{Resumo e o que ainda falta}

Os estudos mostram, portanto, que o resultado do agente piora quando o sistema é maior ou mais complexo, que ele pode perder regras que não consegue ver, que pode tomar atalhos para parecer ter acertado, que erra mais em linguagens antigas e pouco documentadas e que o custo cresce com o volume de código que ele lê. Esses são os desafios do Quadro~\ref{tab:desafios}.

Os estudos que lemos têm algo em comum. Parte deles usa \textit{benchmarks}, que são conjuntos de projetos públicos preparados para testar agentes, como o FreshBrew, o MigrationBench e o RepoMod-Bench. Outra parte usa sistemas de empresas, como o ERP de shoppings do estudo de \citeonline{alves2026legacy}, que é grande, tem mais de vinte anos e foi estudado em um ambiente preparado pelos autores. Nos dois casos, o foco é medir o desempenho do agente: quantos testes passou, quantas regras migrou e quanto custou.

Isso deixa uma pergunta menos explorada. O que acontece quando uma pessoa que desenvolve, no seu dia a dia, pede a um agente que atualize parte de um sistema real e de tamanho menor, e depois precisa conferir o que ele fez? Nesse caso, não existe um conjunto de projetos e testes já preparado. O desenvolvedor precisa descobrir as regras do sistema, definir os resultados esperados e decidir se o código novo pode ser aceito. Nos estudos que lemos, encontramos pouco sobre essa situação.

É esse espaço que o estudo de caso deste trabalho ocupa. O Capítulo~4 explica como o estudo foi feito: o sistema, as regras escolhidas e a forma de conferir o resultado. O Capítulo~5 mostra o que foi encontrado e relaciona esses achados com os desafios do Quadro~\ref{tab:desafios}.


\chapter{Metodologia}
\label{cap:metodologia}
Este capítulo descreve como a pesquisa será conduzida. A Seção~\ref{sec:desenho} apresenta o desenho da pesquisa. A Seção~\ref{sec:selecao} descreve a seleção e a extração do trecho de código estudado. A Seção~\ref{sec:pipeline} detalha o \textit{pipeline} de transformação apoiado por LLM. A Seção~\ref{sec:criterios} define os critérios de avaliação.

\section{Desenho da Pesquisa}
\label{sec:desenho}

Este trabalho é uma pesquisa aplicada, de abordagem qualitativa e quantitativa, conduzida como estudo de caso exploratório \cite{gil2002,prodanov2013}. A pesquisa é aplicada porque busca gerar conhecimento para a solução de um problema concreto: a modernização de uma aplicação \textit{low-code} em uso. É exploratória porque o tema, a extração de regras de negócio de fórmulas Power Fx com apoio de LLMs, ainda é pouco estudado, e o objetivo é compreender o fenômeno e levantar evidências iniciais, e não testar hipóteses de forma generalizável.

O estudo de caso foi escolhido por permitir investigar o problema em profundidade, no seu contexto real, quando os limites entre o fenômeno e o contexto não são claramente definidos \cite{gil2002}. O caso estudado é o sistema \textit{Eventos}, uma aplicação desenvolvida em Microsoft Power Apps e utilizada pela empresa em que a pesquisa é realizada. A unidade de análise é um trecho de lógica de negócio escrito em Power Fx, extraído desse sistema conforme os critérios da Seção~\ref{sec:selecao}.

A pesquisa segue quatro etapas: seleção e extração do trecho a ser estudado; aplicação do \textit{pipeline} de transformação via LLM; execução dos testes de equivalência; e validação qualitativa com \textit{stakeholders} da empresa. Os resultados de cada etapa serão apresentados no Capítulo~\ref{cap:resultados}.

\section{Seleção e Extração da Cobaia}
\label{sec:selecao}

Chama-se aqui de \textit{cobaia} o trecho de código Power Fx que será submetido ao processo de modernização. Como o sistema \textit{Eventos} possui muitas telas e fórmulas, não é viável, no escopo deste trabalho, migrá-lo integralmente. Por isso, será selecionado um trecho representativo do problema descrito na Seção~\ref{sec:lowcode}, de acordo com os seguintes critérios:

\begin{itemize}
    \item \textbf{relevância para o negócio:} o trecho deve implementar uma regra que os usuários do sistema reconheçam como importante, como validações, cálculos ou decisões de fluxo;
    \item \textbf{complexidade lógica:} o trecho deve conter condições aninhadas ou combinadas, operações sobre fontes de dados e mais de um caminho de execução;
    \item \textbf{acoplamento à interface:} a regra deve depender do estado de componentes visuais, de modo que sua extração exija separar a lógica da apresentação;
    \item \textbf{observabilidade:} deve ser possível executar o trecho no ambiente do Power Apps com entradas controladas, para que seu comportamento sirva de referência nos testes.
\end{itemize}

A escolha será feita em conjunto com os responsáveis pelo sistema. O trecho selecionado será exportado da aplicação, junto com as informações necessárias para compreendê-lo: os nomes e tipos dos componentes referenciados, a estrutura das fontes de dados utilizadas e as variáveis globais ou de contexto envolvidas.

Antes de qualquer envio a um serviço de LLM, o material passará por anonimização. Nomes de tabelas, colunas, componentes e telas que revelem informações da empresa serão substituídos por identificadores genéricos; dados pessoais, endereços de e-mail, URLs e identificadores internos serão removidos ou substituídos por valores fictícios. A substituição será feita de forma consistente, preservando a estrutura e os tipos das expressões, para que a lógica permaneça intacta. A tabela de correspondência entre os nomes originais e os anonimizados será mantida apenas localmente e não será enviada ao modelo nem publicada neste trabalho. Esse cuidado atende às políticas de segurança da informação da empresa e à Lei Geral de Proteção de Dados Pessoais (Lei nº 13.709/2018).

\section{Pipeline de Transformação via IA}
\label{sec:pipeline}

A transformação do trecho selecionado será feita por um \textit{pipeline} de engenharia de \textit{prompt} em três etapas sequenciais. A divisão em etapas segue a ideia de passos pequenos e verificáveis defendida por \citeonline{fowler2018refactoring}: em vez de pedir ao modelo uma conversão direta, cada etapa produz um artefato intermediário que pode ser revisado antes de servir de entrada para a etapa seguinte. A Figura~\ref{fig:pipeline} resume o fluxo.

\begin{figure}[H]
    \centering
    \caption{\textit{Pipeline} de transformação via LLM.}
    \label{fig:pipeline}
    \small
    \begin{tabular}{c}
        \fbox{\parbox{0.6\textwidth}{\centering Trecho Power Fx anonimizado + contexto}} \\
        $\downarrow$ \\
        \fbox{\parbox{0.6\textwidth}{\centering \textbf{Etapa 1:} Documentação Semântica}} \\
        $\downarrow$ \textit{revisão humana} \\
        \fbox{\parbox{0.6\textwidth}{\centering \textbf{Etapa 2:} Tradução para Python}} \\
        $\downarrow$ \\
        \fbox{\parbox{0.6\textwidth}{\centering \textbf{Etapa 3:} Geração de Testes em Pytest}} \\
        $\downarrow$ \\
        \fbox{\parbox{0.6\textwidth}{\centering Avaliação (Seção~\ref{sec:criterios})}} \\
    \end{tabular}
    \legend{Fonte: elaboração da autora.}
\end{figure}

\subsection{Etapa 1: Documentação Semântica}

Na primeira etapa, o modelo recebe o trecho Power Fx anonimizado e o seu contexto (componentes, fontes de dados e variáveis) e é instruído a descrever, em linguagem natural e em formato estruturado, as regras de negócio implementadas. A descrição deve identificar as entradas de cada regra, as condições avaliadas, os resultados possíveis e os efeitos colaterais, como gravação de dados, notificações e navegação entre telas.

O objetivo é realizar a engenharia reversa semântica discutida por \citeonline{pereira2026modernizacao}: recuperar o que a fórmula faz, separado de como ela está escrita. O documento gerado será revisado pela autora, comparando cada regra descrita com o código original. Regras omitidas, incorretas ou inventadas pelo modelo serão registradas, pois indicam a qualidade da etapa, e corrigidas antes de seguir. O documento revisado passa a ser a especificação de referência para as etapas seguintes.

\subsection{Etapa 2: Tradução para Python}

Na segunda etapa, o modelo recebe o trecho original e a documentação revisada e é instruído a implementar as regras em Python, respeitando as seguintes restrições:

\begin{itemize}
    \item as regras de negócio devem ser implementadas como funções que não dependam de componentes visuais, recebendo como parâmetros os valores que, no Power Apps, vinham da interface;
    \item o acesso a dados deve ficar separado das regras, por meio de funções ou classes específicas, de modo que as regras possam ser executadas sem conexão com a fonte de dados;
    \item as reações de interface (\texttt{Notify}, \texttt{Navigate}) devem ser substituídas por valores de retorno que indiquem o resultado da decisão;
    \item o código deve usar anotações de tipo e nomes coerentes com a documentação da Etapa 1.
\end{itemize}

O resultado é um \textit{backend} desacoplado, no sentido discutido na Seção~\ref{sec:refatoracao}: a lógica fica isolada e pode ser executada e testada sem a interface.

\subsection{Etapa 3: Geração de Testes em Pytest}

Na terceira etapa, o modelo recebe a documentação revisada e o código Python e é instruído a gerar testes unitários com o \textit{framework} Pytest. Cada regra documentada deve ter ao menos um teste para cada caminho de execução, incluindo valores-limite (por exemplo, a capacidade exatamente atingida) e entradas em branco ou inválidas, que concentram as diferenças de semântica entre Power Fx e Python apontadas na Seção~\ref{sec:ia}.

Há, nessa etapa, um risco de circularidade: se o mesmo modelo gera o código e os resultados esperados dos testes, um erro de interpretação pode aparecer em ambos e os testes passariam mesmo com o comportamento incorreto. Para evitar isso, os testes gerados pelo modelo serão usados apenas para definir os \textit{casos} (as combinações de entrada). Os resultados esperados serão obtidos executando-se o trecho original no Power Apps, em ambiente de teste, com as mesmas entradas. Dessa forma, o sistema original funciona como oráculo, à semelhança dos testes de caracterização propostos por \citeonline{feathers2004working}.

\subsection{Registro do experimento}

Para permitir a reprodução e a análise do processo, serão registrados: o modelo de linguagem utilizado e sua versão; os parâmetros de geração; a data de execução; o texto integral de cada \textit{prompt}; as respostas obtidas; e as intervenções manuais realizadas entre as etapas. Os \textit{prompts} serão apresentados em apêndice.

\section{Critérios de Avaliação}
\label{sec:criterios}

A avaliação combina uma dimensão quantitativa, voltada à equivalência funcional, e uma dimensão qualitativa, voltada à utilidade do resultado para a empresa.

\subsection{Teste de equivalência funcional}

A implementação em Python será considerada equivalente ao trecho original, para os casos avaliados, quando produzir o mesmo resultado que o Power Apps para todas as entradas do conjunto de testes. Esse critério segue o adotado por \citeonline{pan2024lost}: código que apenas executa sem erro não é suficiente; ele precisa reproduzir o comportamento do original. Serão analisados os seguintes indicadores:

\begin{itemize}
    \item \textbf{taxa de aprovação:} proporção de casos de teste em que o resultado do Python coincide com o do Power Apps;
    \item \textbf{cobertura de regras:} proporção das regras presentes no código original que foram identificadas na documentação da Etapa 1;
    \item \textbf{cobertura de código:} proporção das linhas e ramificações do código Python exercitadas pelos testes;
    \item \textbf{intervenções manuais:} número e natureza das correções feitas pela autora em cada etapa do \textit{pipeline}.
\end{itemize}

Cada divergência encontrada será analisada para identificar sua origem (erro de interpretação na documentação, erro de tradução ou diferença de semântica entre as linguagens). Ressalta-se que os testes mostram a equivalência apenas para os casos executados e não provam a equivalência para todas as entradas possíveis; por isso, a escolha dos casos, com atenção a caminhos de execução e valores-limite, é parte do método.

\subsection{Validação qualitativa com \textit{stakeholders}}

A segunda dimensão avalia se a documentação e o código gerados são úteis para quem conhece e mantém o sistema. Serão realizadas entrevistas semiestruturadas com \textit{stakeholders} da empresa, como os responsáveis pelo sistema \textit{Eventos} e usuários que conhecem as regras de negócio envolvidas. Os participantes analisarão a documentação gerada e a estrutura do código Python e responderão a questões sobre:

\begin{itemize}
    \item a fidelidade da documentação às regras de negócio que conhecem;
    \item a clareza da documentação em comparação com a leitura direta das fórmulas Power Fx;
    \item a percepção sobre a manutenibilidade da versão em Python;
    \item a viabilidade de adotar o processo na modernização de outras partes do sistema.
\end{itemize}

As respostas serão analisadas de forma qualitativa, buscando pontos de concordância, divergências e sugestões de melhoria.

\section{Limitações}

O estudo se limita a um único sistema e a um trecho selecionado dele, e seus resultados não podem ser generalizados para todas as aplicações Power Apps nem para outras plataformas \textit{low-code}. A migração da interface de usuário e a integração com fontes de dados de produção estão fora do escopo. O comportamento de LLMs pode variar entre versões e execuções, o que é mitigado, mas não eliminado, pelo registro detalhado descrito na Seção~\ref{sec:pipeline}. Por fim, a validação qualitativa depende de um número reduzido de participantes, todos vinculados à mesma organização.


\chapter{Resultados e Considerações Finais}
\label{cap:resultados}
% TODO: este capítulo será escrito depois do estudo de caso.
% TODO: depois do estudo de caso, escrever aqui o que foi encontrado:
% quais desafios apareceram, quais não apareceram e o que o desenvolvedor pode fazer.
% Roteiro (escrever depois do estudo de caso):
% - os casos de teste e os dados sintéticos usados;
% - o resultado de cada funcionalidade atualizada (os resultados foram iguais aos do sistema atual?);
% - os problemas encontrados, classificados pelos desafios do Quadro da literatura;
% - o tempo gasto e o uso do agente;
% - o que o desenvolvedor pode fazer para lidar com cada desafio observado;
% - os limites do estudo.


\chapter{Conclusão}
\label{cap:conclusao}
Este trabalho teve como objetivo entender os desafios de atualizar sistemas legados com agentes de IA, do ponto de vista do desenvolvimento.

Na literatura, encontramos desafios em seis pontos: tamanho do sistema, complexidade das funcionalidades, regras escondidas, atalhos do agente, linguagens antigas e custo. Em todos os estudos, o resultado do agente precisou ser verificado, e os testes foram a forma mais usada para isso.

A proposta do estudo de caso mostra como observar esses desafios em um sistema real: usar o agente em partes pequenas, proteger os testes e registrar cada problema encontrado.

% TODO: depois do estudo de caso, escrever aqui o que foi encontrado:
% quais desafios apareceram, quais não apareceram e o que o desenvolvedor pode fazer.

Como trabalhos futuros, o mesmo método pode ser aplicado em outros sistemas, com outros agentes e outras trocas de tecnologia, para comparar os resultados.

    

% --------------------------------------
% REFERÊNCIAS
% --------------------------------------
% Colocar SUAS referências no arquivo TCC.bib
% Pode ser criada com o BibDesk - https://bibdesk.sourceforge.io/
% Ou manualmente seguindo o arquivo que deixei como modelo
% --------------------------------------
\bibliography{TCC}

\end{document}
